\documentclass[10pt]{article}
\usepackage{epsf}
\usepackage{amsmath}
\usepackage{amssymb}
\usepackage{palatino}
\usepackage[pdftex]{graphics}
\usepackage{fancyhdr}
\usepackage{epsfig}
\usepackage{multirow}
\usepackage{multicol}
\usepackage{cancel}
\usepackage{hyperref}
\usepackage{longtable}
\usepackage{ulem}
\usepackage[yyyymmdd]{datetime}
\renewcommand{\dateseparator}{-}
\usepackage{listings}

\parindent 0in
\parskip 1ex
\oddsidemargin  0in
\evensidemargin 0in
\textheight 8.5in
\textwidth 6.5in
\topmargin -0.25in

\pagestyle{fancy}
\lhead{\textbf{\leftheader}}
\rhead{\textbf{\rightheader}}
\lfoot{Compiled: \today\ \currenttime}
\cfoot{}
\rfoot{\thepage}

\newcommand{\leftheader}{BME290L: MedTech Prototyping Skills}
\newcommand{\rightheader}{Palmeri (Spring 2023)}

\begin{document}

\title{Final Technical Report} 
\author{MedTech Prototyping Skills (BME290L)}
\maketitle

\textbf{Due Date:} April 29, 2023 @ 09:00\\
\textbf{Late Due Date (no penalty):} May 05, 2023 @ 09:00\\

This final technical report is the culmination of almost all of the skills that
you have developed in this course.  Technical documentation is just as important
in engineering as being savvy at CAD/ECAD.  This report will be a single PDF
file that you will upload to the associated Gradescope assignment, due the last
week of classes.

The following sections should be included in your report\dots

\section{Needs / Specifications}
Summarize the stated user needs and associated specifications for the electronics and the enclosure.
\subsection{Enclosure}
\begin{itemize}
    \item Ergonomically held in one hand
    \item Ability to “easily” push the button to turn on the blinking lights with the hand holding the device
    \item Can be stable placed on a flat surface and operated without picking it up from that flat surface.
    \item Batteries can be “easily” replaced.
    \item Power switch will not be inadvertently switched ”off” during operation.
    \item Potentiometer knobs are accessible, but can use a second hand to adjust them.
    \item Board is securely mounted using your 4 mounting holes.
    \item Test pads will be accessible for testing the electronics when the board is mounted inside (assuming
    a cover of some sort can be removed to access the board). In other words, the entire board does not
    need to be removed for any debugging / testing.
    \item Wires between the PCB connectors and the peripheral components are purposely routed / secured
    internally.
    \item All peripheral components securely mounted in outer shell of the enclosure.
    \item Batteries securely mounted using 2x2 or 1x4 AA battery holders.
    \item Size (smaller is better)
    \item Weight (again, smaller is better)
    \item Weight balance (”hand feel”)
    \item Ergonomics / aesthetics. Try to avoid corners and straight edges. Curves are good (and not sharp).
    \item Stress concentrations if your device fell from 6 ft.
\end{itemize}

\subsection{Electronics}
\begin{itemize}
    \item Specify nominal function of all external items the user interacts with
    (switches, buttons, knobs, LEDs, etc.)
    \item Assign specifications (testable target values) for the LED functions
    (e.g., astable frequency, monostable period).  We will assume that "passing"
    is having a 95\% confidence interval\footnote{For normally-distributed data,
    95\% CI $\approx$ 2 SD.} that falls within $\pm$10\% of these nominal
    values.
\end{itemize}


\section{Design Documents}

\subsection{Functional Decomposition}
Please include a complete functional decomposition of your device.
\subsection{ECAD}
\subsubsection{Schematic Capture}
Be sure to include all sheets.  Demonstrate all best practices in your
schematics.

\subsubsection{PCB Layout}
Show the front and back layers, separately.  Instead of "plotting" the layers,
you can take screenshots that best show each layer, with the other layer
"hidden".

\subsubsection{Bill of Materials}
This should be generated using KiCad's tool, including footprints and items
excluded from the PCB.  Consider using the
\href{https://github.com/openscopeproject/InteractiveHtmlBom}{Interactive HTML
BOM} extension.

\subsection{CAD}
\begin{itemize}
    \item Orthographic mechanical drawings of all parts.
    \item 3D renderings of the exploded assembly.
\end{itemize}

\section{Testing \& Analysis}

\textbf{Note: If your PCB is not working in a reasonable timeframe, we will find
another PCB for you to perform all of your testing on.}
\subsection{Testing to Specification}
For each user need / specification outlined in the first section, test how well
you met each need / specification, assuming that your test passes if your 95\%
CIs are within $\pm$10\% of the nominal specifications.  For measurable
quantities with inherent variability, this means that you need to have replicate
measurements of these values!

Oscilloscope output can be captured using cellphone screenshots with
measurements being displayed on the scope.

For any tests that do not pass, please provide an explanation for why the test
didn't pass and what you might change in your design to increase the likelihood
of passing if your design was revised.  For example:

\begin{quote}
    The nominal astable frequency was 10$\pm$1 Hz, but my astable output was
    9.1$\pm$0.5 Hz.  This performance fails at 95\% CI.  One reason for this
    failure was associated with the 15\% tolerance on the resistors and
    capacitors used to configure this astable oscillator.  As shown in the table
    below, Rx and Ry varied by over 10\% from the nominal values, leading to a
    negative bias in the output frequency.
\end{quote}

\subsection{Test Power Consumption}
In addition to the specifications that you designed your device around, please
provide an analysis of the battery life of your device.  This can be estimated
by measuring the current demands of your circuit when "idle" and when blinking
the LED.

You will want to indirectly estimate this battery life, and not try to test this
by simply running your device with batteries until they fully discharge.

\subsection{Optimizations}
For the specifications that could be optimized (e.g., size, weight, etc.) comment
on how your design optimized these quantities, and given your current insights,
how your design could be revised to improve upon these criteria.

\section{Microcontroller}
You will be uploading your microcontroller firmware as a zip archive to a
separate Gradescope assignment, using all of the best practices discussed in
lecture.  These include:
\begin{itemize}
    \item Organized inclusion of libraries
    \item Definition of preprocessor macros
    \item Consolidated initialization of global variables
    \item Source code without "hard-coded" functional values
    \item Logical, "action-based" main loop
\end{itemize}

In this report, include:
\begin{enumerate}
    \item Testing data of the same LED outputs as your did for the discrete
    electronics version of your device.
    \item Please include a discussion of the "pros" and "cons" of implementing
    the desired LED functions using discrete electronics versus a
    microcontroller.
\end{enumerate} 

\end{document}